\begin{lemma}[Proved: two flag-complex budgets]
\label{CP2:A-high-flag}
Let $G$ be a finite connected graph of cycle rank $R$, and let $X(G)$
be its clique complex.  If $R\leq6$, then
$H_2(X(G);\mathbb Q)=0$.  If $R\leq8$, then
$H_4(X(G);\mathbb Q)=0$; in this latter range $G$ contains at most
one clique of size five and none of size six.
\end{lemma}
\begin{proof}
Suppose first that second homology is nonzero.  Among induced
subgraphs having this property choose one with the fewest vertices,
and pass to its relevant connected component.  Deleting any vertex
$v$ kills second homology.  The deletion together with the star of
$v$ covers the clique complex, and their intersection is the link
of $v$.  Mayer--Vietoris therefore injects the nonzero second
homology into the first homology of that link.  A clique complex
on at most three vertices has zero first homology, so every vertex
has at least four neighbors.  There are at least six vertices:
with five vertices the degree condition forces a complete graph,
whose clique complex is contractible.  If their number is $a$ and
the edge number is $b$, then $b\geq2a$ and the cycle rank is at
least $a+1\geq7$.  Cycle rank is monotone under taking subgraphs,
which proves the first assertion.

Two distinct five-cliques sharing $t$ vertices have union with
cycle rank at least $12-\binom t2+t-1$ for $1\leq t\leq4$;
its minimum is nine, at $t=4$.  Disjoint five-cliques contribute
twelve independent cycles.  A six-clique already has cycle rank
ten.  Thus at most one four-simplex occurs when $R\leq8$, and
there is no five-simplex.  Its boundary, if it occurs, is a nonzero
sum of its five distinct faces.  The fourth chain boundary is
therefore injective, proving the second assertion.
\end{proof}

\begin{lemma}[Proved: a graph cofiber over the five-element band]
\label{CP2:A-high-P5-graph}
Let $P=\{1\}\sqcup\{x_{ij}:0\leq i,j\leq1\}$, with
$x_{ij}x_{kl}=x_{il}$.  Let a finite graph $\Gamma$ carry a
right $P$ action sending each oriented edge either to an
oriented edge by its linear parameter or to a vertex.
The cofiber of $\Gamma_{hP}\to BP$ has zero rational
degree-two products.  Consequently this conclusion holds
for every collage between $P$ and
$R_0=\{1,p,q\}$, $ab=b$ for $a,b\in\{p,q\}$,
with any nonempty finite $(R_0,P)$-biset.
\end{lemma}
\begin{proof}
Write $A=\mathbb QP$, $e=x_{00}$, $f=1-e$,
$a=x_{10}-x_{00}$ and $b=x_{01}-x_{00}$.  There is an exact
projective resolution of the trivial left module
\[
 0\longrightarrow Ae\xrightarrow{\cdot b}Af
   \xrightarrow{\cdot a}Ae\longrightarrow\mathbb Q
   \longrightarrow0.
\]
To check it directly, bases for the first and last projective
are $x_{00},x_{10}$; a basis for the middle projective is
$f,b,c=x_{11}-x_{10}$.  The first map sends the basis to
$b,c$, and the second sends $f$ to $a$ and kills $b,c$.
The last augmentation has kernel $\mathbb Qa$.

The images of the two actors in column $j$ are the same
retraction subgraph $\Gamma_j$.  The maps $x_{01}$ and
$x_{00}$ restrict to inverse isomorphisms between
$\Gamma_0$ and $\Gamma_1$.  If these subgraphs share a
vertex, it is fixed by every band element, hence by $P$.
It supplies a section of $\Gamma_{hP}\to BP$.
Mayer--Vietoris then makes all second classes of the
cofiber relative to $BP$.  They lift from its suspension
quotient and have zero products.

Otherwise the two subgraphs have disjoint vertex and edge
sets.  Put $W=H_1(\Gamma;\mathbb Q)$.  Tensoring the
displayed resolution with $W$ identifies
\[
 \operatorname{Tor}_2^A(W,\mathbb Q)
   =\ker\bigl(We\xrightarrow{\cdot b}Wf\bigr).
\]
Here $We=H_1(\Gamma_0;\mathbb Q)$, and $wb$ is the
difference between the image of the cycle $w$ in
$\Gamma_1$ and its original cycle in $\Gamma_0$.
A graph has no two-dimensional boundaries, and these
two cycles have disjoint edge supports.  Their difference
is zero only if $w=0$.  Thus this Tor group is zero.

This also proves the required fourth-homology vanishing
when $\Gamma$ is disconnected.  In fact, put
$L=\widetilde H_0(\Gamma;\mathbb Q)$ and let $K$ be the
two-term complex
$[C_1\Gamma\to\ker(C_0\Gamma\to\mathbb Q)]$.
The triangle $W[1]\to K\to L$ and the length-two
resolution show that $H_3(K\otimes_A^{\mathbf L}\mathbb Q)
=\operatorname{Tor}_2^A(W,\mathbb Q)$: the possible
outer terms are $\operatorname{Tor}_4^A(L,\mathbb Q)$
and $\operatorname{Tor}_3^A(L,\mathbb Q)$, both zero.
The augmented action-bar cofiber has this group as its
fourth homology.  Therefore its degree-two products are
zero in the disjoint-image case also.

For the biset conclusion, the left $R_0$-action consists
of two retractions $P_0,Q_0$ with a common image $U$.
Its action category has a natural graph model: vertices
are $U$, and each state $t\notin U$ gives an edge from
$P_0t$ to $Q_0t$.  Indeed, a state outside $U$ has only
its identity and its two outgoing arrows to $U$.  The
endomorphism monoid of each vertex of $U$ is $R_0$,
whose classifying space is contractible.  The two
outgoing-arrow cones attach along these contractible
spaces; collapsing them gives the graph with the outside
state as a bivalent vertex, then its stated edge.
On nerve simplices this collapse uses the sum of the
barycentric coordinates after the initial outgoing
arrow, so it commutes with every equivariant set map.
Such a map sends an edge to the edge of the image state
or collapses it when that state lies in $U$.  The
commuting right $P$ action therefore gives exactly the
graph action already treated.  Since $BR_0$ is
contractible, the collage is the stated cofiber.
\end{proof}

\begin{corollary}[Proved: two-object acyclic support through eleven arrows]
\label{CP2:A-high-two-objects}
A connected category with two objects, acyclic support, and at most
eleven arrows cannot have the homotopy type of $\mathbb{CP}^2$.
More strongly, its rational degree-two products are zero.
\end{corollary}
\begin{proof}
Write $a,b$ for the two corner orders and $t\geq1$ for
the cross-hom size.  If $t=1$, the collage is the join
of its two connected corner spaces and has zero second
cohomology.  If $t=2$, use
Lemma~\ref{CP2:A-high-two-cross}.  Suppose $t\geq3$.
If both corners have zero second cohomology, the
two-block lemma applies.  Otherwise one corner has
order at least five by Corollary~\ref{CP2:A-five-types}.
Its order is at most seven and the other corner has
order at most three.  With other corner of order at most
two, exact orbit comparison reduces to
Corollary~\ref{CP2:A-high-seven-act}.
The sole remaining budget is $(a,b,t)=(5,3,3)$ or its
opposite, and the positive corner is $P$ of the preceding
lemma.  If the order-three corner has a projective trivial
module on the required side, orbit comparison gives a
two-row act cone.  Otherwise its minimal ideal is the
two-element right-zero band, after choosing the direction
of the collage, and Lemma~\ref{CP2:A-high-P5-graph}
applies.  All connected cases have zero products.
\end{proof}

\begin{lemma}[Proved: a cut containing all parallel hom-sets]
\label{CP2:A-high-parallel-cut}
Let $C$ have acyclic support of cycle rank at most six.
Suppose no two hom-sets of size greater than one occur
serially: for such $C(x,y),C(z,w)$ one never has $y\leq z$.
If every global degree-two class restricts to zero at
each corner with nonzero second cohomology, all rational
degree-two products on $BC$ vanish.
\end{lemma}
\begin{proof}
We use the following join observation, including disconnected
predecessor categories.  If $D$ is nonempty and $M$ is a finite
monoid, then $BD*BM$ has zero rational second cohomology.
Indeed $BM$ is connected and $H_1(BM;\mathbb Q)=0$, since
its universal group is finite.  The reduced join Kunneth
formula has no nonzero summand in degree two.  The inclusion
of $BM$ in this join is null-homotopic, by choosing any point
of $BD$.  These facts also justify the join cuts below.

Let $D$ be the downset generated by the sources of the
parallel hom-sets.  Each such source lies in $D$, and
its target cannot lie there: otherwise that target is
at or below the source of another parallel hom-set,
contrary to the assumption.  Every parallel hom-set
therefore crosses from $D$ to its complementary upset.
The full categories on either side have singleton cross
hom-sets.  Each of their components with more than one
object has second cohomology equal to that of its support
by Lemma~\ref{CP2:A-singleton-support}, and the latter
vanishes by Lemma~\ref{CP2:A-high-flag}.  On a one-object
component, global classes restrict to zero by the
hypothesis.  All global classes therefore vanish on
both blocks.  Lemma~\ref{CP2:A-two-block} proves the claim.
When there are no parallel hom-sets, apply the same
singleton-support argument directly, component by component.
\end{proof}

\begin{lemma}[Proved: two parallel excesses]
\label{CP2:A-high-two-parallel}
Suppose $C$ is connected with acyclic support, every
corner has order at most four, its support cycle rank
is at most six, and
$\sum_{x<y}(|C(x,y)|-1)\leq2$.  If
$\dim H^2(BC;\mathbb Q)\leq1$, its degree-two squares vanish.
\end{lemma}
\begin{proof}
All corners are rationally acyclic.  Unless two doubled
hom-sets occur serially, the preceding lemma applies.
Write the two serial hom-sets as $a\to b$ and $c\to d$,
where $b\leq c$.  If $b<c$, the full downset of $c$ has
singleton hom-sets into $c$ and is the categorical join
of its predecessor category and its corner.  This join
has zero second rational cohomology: the corner is
rationally acyclic.  Its complement has only singleton
cross hom-sets and zero second cohomology by the support
budget.  The two-block lemma applies.

It remains that the two doubled sets are $a\to b$ and
$b\to c$.  If $b<x$ with $x\ne c$, the full downset
$D$ of $b$ has only singleton arrows into $x$.
Consequently $D\cup\{x\}$ is a join with the rationally
acyclic corner at $x$ and has zero second cohomology.
Every global second class vanishes on $D$.  The complement
has singleton cross hom-sets and zero second cohomology,
so the two-block lemma again applies.  An extra predecessor
of $b$ is handled by opposition.

Otherwise $b$ has just the two neighbors $a,c$.  The
category is the actual nerve union of the full triangle
on $a,b,c$ and the deletion of $b$.  Their intersection
is the singleton edge $a\to c$, a rationally acyclic
join.  The deletion has zero second cohomology by the
singleton-support lemma.  Mayer--Vietoris identifies
the triangle's second cohomology with the global group
and makes fourth-cohomology restriction to the two
pieces injective.  Its hom sizes are $(2,2,1)$, so
Lemma~\ref{CP2:A-high-small-binary} kills its rank-one
square.  Its restriction to the other piece is already
zero, proving the result.
\end{proof}

\begin{lemma}[Proved: finite affine acts and two states per object]
\label{CP2:A-high-affine}
Let $T\subset\mathbb Q$ be a nonempty finite right monoid act,
whose maps are affine of slopes $0,1,-1$.  Adjoin a trivial sink
to the monoid with outgoing arrows $T$.  Every positive rational
cohomology product of this act-cone category is zero.

More generally, the same conclusion holds for an act cone over
any finite category $S$ when its right act $F$ has one or two
states at each object.  Both statements have left-act duals.
\end{lemma}
\begin{proof}
A slope-one affine self-map of a finite set with at least two
elements is the identity.  A slope-minus-one map is its unique
reflection about the midpoint of its extreme elements.  Extend
the action to the interval $I=[\min T,\max T]$.  All affine
identities holding on $T$ hold on $I$, so this is a strict action.
The interval bar maps equivalently to the monoid classifying
space, level by level, because $I$ is contractible.  Consequently
the act cone is equivalent to the cofiber of the inclusion of
the $T$ bar in the $I$ bar.  This is the pointed bar with fiber
$I/T$, with its basepoint bar collapsed.  The fiber is a wedge
of circles indexed by consecutive elements of $T$.  Its actors
are identities, a reflection permuting and reversing circles,
or pointed constants.  On each circle take the map of degree two.
It commutes strictly with all these actors: in a circle coordinate,
$2(-t)=-2t$ modulo one.  It induces a self-map of the pointed bar.
The reduced cellular chain complex of the fiber is concentrated
in degree one, and this self-map is multiplication by two on
that whole complex.  The same is true in every degree of the
reduced bar total complex, hence on all its positive cohomology.
For positive classes $u,v$, functoriality gives $2uv=4uv$, and
thus $uv=0$ over $\mathbb Q$.  A singleton act has contractible
cone and presents no exception.

For the categorical assertion, replace each $F(s)$ by its convex
simplex $I(s)$ and each set map by its affine extension.  These
extensions compose exactly.  Each $I(s)$ is a point or interval,
and each quotient $I(s)/F(s)$ is a point or circle.  The induced
maps are pointed constants or circle maps of degree $1$ or $-1$.
Their degree-two self-maps form a strict natural transformation.
The cofiber of $\operatorname{hocolim}F\to\operatorname{hocolim}I$
is the pointed bar of these quotients, and
$\operatorname{hocolim}I\to BS$ is an equivalence.  It is thus
the required act cone.  The identical reduced-chain argument
proves the claim.  Reversing arrows proves the dual statements.
\end{proof}

\begin{lemma}[Proved: two minimal rows]
\label{CP2:A-high-two-row-act}
Suppose the minimal ideal of a finite monoid $M$ has at most
two rows.  For every nonempty finite right $M$-act $T$, all
rational degree-two products of its act cone vanish.  The
left-act statement holds with columns in place of rows.
\end{lemma}
\begin{proof}
Write $R=\mathbb QM$.  For a minimal idempotent $e$, the corner
$H=eMe$ is a finite group.  Average it to the algebra idempotent
$f=|H|^{-1}\sum_{h\in H}h$.  The projective left module $Rf$
has one basis vector per minimal row: the right $H$-orbits of
$Me$ are free, and multiplying by $f$ averages each such orbit.
Left multiplication acts by maps of the row set.  The row
augmentation $Rf\to\mathbb Q$ is onto.  With one row it is an
isomorphism, so the trivial left module is projective.
With two rows its kernel is a one-dimensional module
$\mathbb Q_\chi$, with $\chi(m)\in\{0,1,-1\}$ and
$\chi(mn)=\chi(m)\chi(n)$.

Put $L_T=\ker(\mathbb QT\to\mathbb Q)$ and
$D L_T=\operatorname{Hom}_{\mathbb Q}(L_T,\mathbb Q)$.
The augmented action bar gives, in degrees $j\geq2$,
\[
 H_j(\operatorname{cone}(T_{hM}\to BM);\mathbb Q)
       =\operatorname{Tor}_{j-1}^{R}(L_T,\mathbb Q).
\]
Indeed, the cone chain complex is the bar of the two-term
complex $[\mathbb QT\to\mathbb Q]$, which is
quasi-isomorphic to $L_T[1]$.  Dualizing a projective resolution
identifies degree-two cohomology with
$\operatorname{Ext}^1_R(\mathbb Q,D L_T)$.  The row augmentation
sequence makes the connecting map
\[
 \operatorname{Hom}_R(\mathbb Q_\chi,D L_T)
       \longrightarrow H^2(\operatorname{cone}(T_{hM}\to BM);\mathbb Q)
\]
surjective.  These constructions are natural in maps of acts.

Represent a class by a linear functional $a:L_T\to\mathbb Q$
satisfying $a(zm)=\chi(m)a(z)$.  Choose $t_0\in T$ and set
$q(t)=a(t-t_0)$.  Then
\[
 q(tm)=\chi(m)q(t)+q(t_0m).
\]
Thus $q$ is an equivariant map to the finite affine act $q(T)$.
On its augmentation kernel, the functional sending a formal
sum $\sum b_xx$ to $\sum b_xx\in\mathbb Q$ pulls back to $a$.
Naturality of the connecting map shows that the original class
pulls back from this affine act cone.  Its square vanishes by
Lemma~\ref{CP2:A-high-affine}.  Every square therefore vanishes,
and rational polarization kills all mixed products.  The one-row
case has zero second cohomology by projectivity.  Opposition
interchanges rows and columns.
\end{proof}

\begin{corollary}[Proved: act cones through order seven]
\label{CP2:A-high-seven-act}
Every act cone over a monoid of order at most seven has zero
rational degree-two products.  There is no bound on the number
of states and no bounded-homology hypothesis.
\end{corollary}
\begin{proof}
For a right act use the preceding lemma if the minimal row
number is at most two.  If the column number is one, the
group-corner comparison of Lemma~\ref{CP2:A-small-minimal-ideal}
gives $BM\simeq BG$ for a finite group, so $H^2(BM;\mathbb Q)=0$.
All second cohomology of the cone then lifts from its suspension
quotient, whose positive products vanish.  The only other
possibility at order at most seven has three rows, two columns,
trivial minimal group, and just the global identity outside the
minimal ideal.  This is the identity-adjoined $3\times2$ band.
The projective resolution in Lemma~\ref{CP2:A-band-spheres}
has length two over $\mathbb Q$ and resolves the trivial module.
Thus its act cone has
$H_4=\operatorname{Tor}_3^R(L_T,\mathbb Q)=0$.
This exhausts the row, column and group-size count.  The left-act
case follows by opposition.
\end{proof}

\begin{lemma}[Proved: exact orbit comparison]
\label{CP2:A-high-orbits}
In a collage with one monoid endpoint $N$, suppose its trivial
rational module is projective on the side used for the
coinvariants.  Replace $N$ by a trivial object and every
incident stalk by its actual $N$-orbit set.  The quotient functor
is a rational homology equivalence and induces an isomorphism
of rational cohomology rings.  The same operation can be made
at an extreme object of an acyclic-support category.
\end{lemma}
\begin{proof}
There is an actual quotient category: endpoint actions identify
only arrows incident with that endpoint, and compatibility of
composition with the actions defines all quotient products.
The action bar of an $N$-set maps to its discrete orbit set.
Its rational homology is the derived coinvariant functor on
the permutation module.  Projectivity of the trivial module
makes that functor exact, so this map is a rational homology
equivalence.  The vertex map $BN\to *$ is one also.
In each degree of every other bar direction the same argument
applies.  Filtering the resulting first-quadrant double bars
proves a rational homology equivalence on the cross-arrow
spaces.  The two-block homotopy pushout of
Lemma~\ref{CP2:A-two-block} completes the comparison.
For an extreme object with several incident stalks, the same
argument is natural in their remaining category variable.
As the comparison is induced by a functor, its cohomology
isomorphism respects products.
\end{proof}

\begin{lemma}[Proved: a small corner with binary stalks]
\label{CP2:A-high-small-binary}
Let a finite category be a collage of a source block $S$ and
one sink corner $N$ of order at most four.  Suppose every
stalk from an object of $S$ has one or two elements.  If the
second rational cohomology of the whole collage has dimension
at most one, every degree-two square is zero.  There is a
dual assertion for a small source corner.
\end{lemma}
\begin{proof}
If the trivial right $\mathbb QN$-module is projective,
Lemma~\ref{CP2:A-high-orbits} reduces the claim to
Lemma~\ref{CP2:A-high-affine}.
Otherwise the projectivity test
Lemma~\ref{CP2:A-projectivity-test} and the minimal-ideal
cardinality show that the ideal is a right-zero band of two
or three elements.  It has one row and trivial group: two
rows or a nontrivial group would require at least five
monoid elements when there are two columns.  Its elements
are right zeros of the whole monoid, so $BN$ is contractible.

On a set of at most two states, the left actions of this
right-zero band coincide.  Indeed they are idempotent maps
with a common image.  If that image has two elements all
are identities; if it has one element all are that same
constant.  Denote this retraction on each stalk by $E$.
The family is natural in $S$.  For $n\in N$, the equations
$nE=E=En$ follow because $nk=k$ and $kn$ remains in the
minimal ideal.  Its image is consequently a biset retract
with trivial $N$ action.  The core collage is equivalent
to an objectwise one- or two-state act cone, since $BN$
is contractible.  Its positive products vanish.

The inclusion of the core has an actual retraction.  Its
quotient in the double-mapping-cylinder model collapses
both vertex blocks and is a suspension of the quotient
correspondence.  If the core has zero second cohomology,
all second classes lift from this suspension.  Otherwise
its injective pullback spans the assumed space of dimension
at most one.  In either case every square vanishes.
\end{proof}

\begin{lemma}[Proved: two cross arrows, without corner restrictions]
\label{CP2:A-high-two-cross}
A two-object category with arrows in only one direction and
exactly two cross arrows has zero rational degree-two products.
\end{lemma}
\begin{proof}
Write its two-element biset as $T$, and consider the image of
the left corner in the four-element transformation monoid of
$T$.  If both constants occur, a commuting right transformation
fixes both points, so the right action is trivial.  If exactly
one constant occurs, no interchange can occur; both actions
fix its value.  That common fixed state supplies a section of
the correspondence projection to the product of the two corner
spaces.  Their second cohomology injects into the correspondence,
so Mayer--Vietoris makes all global second classes relative
to the vertex blocks.  Their products vanish by
Lemma~\ref{CP2:A-two-block}.

If no constant occurs, the left image is trivial or contains
the interchange.  In the latter case every commuting map is
a permutation.  The correspondence is the two-sheeted covering
of the product of the corner spaces determined by these two
permutation actions.  Its transfer makes pullback injective
over $\mathbb Q$, so the same relative argument applies.
The remaining case, and the first case after opposition, has
a trivial action at one corner.  Its correspondence is
$Z\times BN$, where $Z$ is the ordinary action bar of the
other corner.  A finite monoid has zero rational first
homology, hence
\[
 H^2(Z\times BN;\mathbb Q)
  =H^2(Z;\mathbb Q)\oplus
     H^0(Z;\mathbb Q)\otimes H^2(BN;\mathbb Q).
\]
The diagonal map from $H^2(BN)$ to the second summand is
injective.  In degrees one and two the Mayer--Vietoris
sequences therefore identify pullback from the ordinary
two-state act cone with an isomorphism onto the second
cohomology of the collage.  This is an actual functor map
and preserves products.  Lemma~\ref{CP2:A-high-affine}
annihilates all those products.
\end{proof}

\begin{lemma}[Proved: a nonlocal algebra of dimension at most four]
\label{CP2:A-high-small-Tor}
For a nonlocal finite-dimensional $\mathbb Q$-algebra $A$ with
$\dim A\leq4$ and finite-dimensional modules $W_A,{}_A V$,
nonzero $\operatorname{Tor}_2^A(W,V)$ forces nonzero Tor in
arbitrarily high degrees.  If a finite monoid $N$ has more
than one element, its algebra $\mathbb QN$ is nonlocal.
\end{lemma}
\begin{proof}
Put $J=\operatorname{rad}A$.  Semisimple algebras contribute
no positive Tor.  If the semisimple quotient has a nonrational
field factor and $J\ne0$, nonlocality and the dimension bound
leave only a quadratic field, a rational factor, and a
one-dimensional radical.  That radical cannot carry a unital
quadratic-field action and is supported at the rational
factor.  By lifting orthogonal idempotents, $A$ is a product
of a semisimple algebra and the dual numbers.  A matrix
simple factor of degree greater than one already uses all
four dimensions and is semisimple.

The remaining case has $A/J=\mathbb Q^r$, $r\geq2$, and
$\dim J\leq2$.  If $J^2\ne0$, necessarily $r=2$,
$\dim J=2$, $\dim J^2=1$, and $J^3=0$.  The unique
one-dimensional component of $J/J^2$ must be a loop:
an off-diagonal component cannot multiply with itself.
A lift $x$ has $x^2\ne0$, so
$A=\mathbb Q\times\mathbb Q[x]/(x^3)$.
For $\mathbb Q[x]/(x^m)$, every finite module splits into
modules $A/(x^a)$, and a nonprojective such module has the
two-periodic free resolution with alternating maps
$x^a,x^{m-a}$.  Nonzero degree-two Tor consequently recurs
in every subsequent even degree.  This treats these cases
and the dual numbers.

It remains that $J^2=0$.  Lift the $r$ orthogonal primitive
idempotents.  The algebra is the radical-square-zero quiver
algebra with $r$ vertices and $d=\dim J$ arrows, where
$r+d\leq4$.  First syzygies $\Omega W,\Omega V$ of minimal
projective covers are semisimple.  The inclusion
$\Omega W\to P_W$ tensors to zero with $\Omega V$,
because its image lies in $P_WJ$ and $J\Omega V=0$.
The two projective-cover sequences give
\[
 \operatorname{Tor}_n^A(W,V)
 =\operatorname{Tor}_{n-2}^A(\Omega W,\Omega V)
 \qquad(n\geq2),
\]
including ordinary tensor product at $n=2$.  A nonzero
tensor means that a vertex $v$ occurs in both syzygies.
It must have an entering and a leaving arrow.  With at
most two arrows and $r+d\leq4$, these arrows form a loop
or a two-cycle through $v$.  The resolution of a vertex
simple in a radical-square-zero algebra has one generator
for every directed path starting there.  Repeating that
cycle gives nonzero Tor from the right $v$-simple to the
left $v$-simple in arbitrarily high degrees.  These are
direct summands of the displayed groups, proving the claim.

Finally, a nonunit in a finite monoid has an idempotent
power different from the identity.  Its basis element is
a nontrivial algebra idempotent.  If there is no nonunit,
the monoid is a nontrivial finite group, whose averaging
idempotent is nontrivial.  A local algebra has no such
idempotent.  Thus $\mathbb QN$ is nonlocal.
\end{proof}

\begin{lemma}[Proved: augmentation kernels between two trivial extremes]
\label{CP2:A-high-linear-Tor}
Let $S$ be a finite category between a new trivial source $a$
and a new trivial sink $b$.  Suppose every incoming and
outgoing stalk at $S$ is nonempty, and put $R=\mathbb QS$.
Let $V$ be the left permutation module of the incoming
stalks and $W$ the right permutation module of the outgoing
stalks.  Their augmentations are onto the constant left
and right modules; write $\bar V,\bar W$ for the kernels.
Allow any nonempty finite set of direct arrows $a\to b$
and any compatible composition.  Then
\[
 H_j(BC;\mathbb Q)=\operatorname{Tor}_{j-2}^R(\bar W,\bar V)
 \qquad(j\geq3).
\]
For a three-object category whose extreme corner algebras
are semisimple and whose middle corner has order at most
four, bounded rational homology implies $H_4(BC;\mathbb Q)=0$.
\end{lemma}
\begin{proof}
Quotient the normalized category chains by the full subcategory
on $a,b$.  This is a graph, with no chains above degree one.
The quotient is the total two-sided categorical bar of the
two complexes $[W\to\mathbf1_r]$ and
$[V\to\mathbf1_\ell]$, with permutation terms in degree
one.  The four horizontal parts are respectively the bar
of $S$, its incoming and outgoing action bars, and its
two-sided action bar, shifted by zero, one, one, and two.
The face composing the two outer arrows lands in the
quotiented graph.  Every other face is the usual bar
face or an augmentation.  This gives an equality of
chain descriptions, including empty internal strings.

The two complexes are quasi-isomorphic to
$\bar W[1],\bar V[1]$.  This replacement preserves the
bar quasi-isomorphism because all tensor factors over
the ground field are exact; filtering by internal degree
is a first-quadrant comparison.  Its homology is therefore
the displayed Tor shifted by two.  The graph long exact
sequence proves the assertion in degrees at least three.
For semisimple extreme corners use the actual orbit
functors of Lemma~\ref{CP2:A-high-orbits}.  If the middle
monoid is trivial there is no positive Tor.  Otherwise
Lemma~\ref{CP2:A-high-small-Tor} applies to its algebra.
A nonzero fourth homology group would imply unbounded
nonzero homology, proving the last claim.
\end{proof}
